% ============================================================
%  DROP-IN REPLACEMENT for Theorem \label{thm:expohdLO}
%  (expoHD on LeadingOnes)
%
%  Replaces:  \begin{restatable}{theorem}{expohdLO}\label{thm:expohdLO} ...
%             through the matching \end{proof}
%
%  WHAT TO DELETE from TEVCNewFullPaper.tex:
%    - The old \begin{restatable}...\end{restatable} block for \label{thm:expohdLO}
%    - The old \begin{proof}...\end{proof} block immediately after it
%      (from the orphan line  $n^{-1/(2n)} \ge ...$  through \end{proof})
%
%  WHAT TO INSERT in its place: everything between the ===== markers below.
%
%  Also apply the DOWNSTREAM EDITS and TABLE 2 CHANGES noted at the
%  bottom of this file.
% ============================================================


\begin{restatable}{theorem}{expohdLO}\label{thm:expohdLO}
The {\oneoneIA } using \IPHfcm{} with {\expoHD } optimises
\textsc{LeadingOnes} in $\Theta(n^{5/2}/\log^2 n)$ expected fitness function
evaluations.
\end{restatable}

\begin{proof}
At level~$i$ (i.e., with $i$ leading $1$-bits), the solution has the form
$x = 1^{i}\,0\,b_{i+2}\cdots b_n$ where the tail bits $b_{i+2},\dots,b_n$
are uniformly distributed.
%~\cite{DJW02}.
The Hamming distance to the optimum~$1^n$ is therefore
$H(x,\text{opt}) = 1 + B$ where $B\sim\mathrm{Bin}(n-i-1,\,1/2)$.

\medskip
\noindent\textbf{Concentration of the mutation potential.}
By a Chernoff bound,
\begin{align*}
&\Pr\!\left[\left|B - \tfrac{n-i-1}{2}\right| > n^{2/3}\right]
\;\le\; 2\exp\!\left(-\frac{2n^{4/3}}{n-i-1}\right)
\\\;&\le\; 2e^{-2n^{1/3}}
\;=\; o(n^{-2}).
\end{align*}

A union bound over all $n$ levels and polynomially many iterations shows that,
with overwhelming probability (w.o.p.), the Hamming distance satisfies
\[
H(x,\text{opt})
\;=\; 1 + \tfrac{n-i-1}{2} \pm n^{2/3}
\;=\; \tfrac{n-i+1}{2} \pm n^{2/3}
\]
throughout the entire run.
Setting $h_i := (n-i+1)/(2n)$ and noting that
$n^{\pm n^{-1/3}} = e^{\pm n^{-1/3}\ln n} = 1 \pm o(1)$,
the mutation potential satisfies
\begin{equation}\label{eq:M-expoHD-conc}
M \;=\; n^{H/n}
  \;=\; (1\pm o(1))\;n^{h_i}
  \;=\; (1\pm o(1))\;n^{1/2-(i-1)/(2n)}.
\end{equation}
If the rare event that $|B - (n-i-1)/2| > n^{2/3}$ occurs, we
pessimistically assume $M=n$, which yields polynomial expected runtime and
contributes $o(1)$ to the total expectation.

\medskip
\noindent\textbf{Per-level cost.}
The first bit flipped by the hypermutation is chosen uniformly from
$\{1,\dots,n\}$.  With probability $i/n$ it hits a leading $1$-bit,
destroying the prefix and consuming the full potential of~$M$ evaluations
with no improvement; with probability $1/n$ it hits the leftmost $0$-bit,
yielding an improvement at cost~$1$; with probability $(n-i-1)/n$ it hits
a tail bit, costing $1$~evaluation with no improvement.
By the law of total expectation,
\[
E(f_i)
= \frac{i}{n}\bigl(M+E(f_i)\bigr)
  +\frac{1}{n}\cdot 1
  +\frac{n-i-1}{n}\bigl(1+E(f_i)\bigr),
\]
and solving gives
\begin{equation}\label{eq:Efi-expoHD}
E(f_i) \;=\; i\cdot M + n - i.
\end{equation}
By~\eqref{eq:M-expoHD-conc}, this equals
$(1\pm o(1))\;i\cdot n^{1/2-(i-1)/(2n)} + n - i$ w.o.p.

\medskip
\noindent\textbf{Evaluating the sum.}
We evaluate $S := \sum_{i=0}^{n-1} E(f_i)$, which governs both bounds.
Factoring out $n^{1/2+1/(2n)}$ from the leading term and setting
$y:=n^{-1/(2n)}<1$,
\begin{align*}
S
&= (1\pm o(1))\sum_{i=0}^{n-1} i\cdot n^{1/2-(i-1)/(2n)}
   \;+\;\underbrace{\sum_{i=0}^{n-1}(n-i)}_{\Theta(n^2)} \\
&= (1\pm o(1))\;n^{1/2+1/(2n)}\sum_{i=1}^{n-1} i\,y^{i}
   \;+\;\Theta(n^2).
\end{align*}
By the standard identity $\sum_{i=1}^{\infty} i\,y^i = y/(1-y)^2$
and Lemma~\ref{lem:expo} (with $\epsilon=1/2$),
\[
1-y \;=\; 1 - n^{-1/(2n)}
    \;=\; \frac{\ln n}{2n}\bigl(1+o(1)\bigr),
\]
so
\[
\frac{y}{(1-y)^2}
\;=\; \frac{1+o(1)}{\bigl(\frac{\ln n}{2n}\bigr)^2}
\;=\; \frac{4n^2(1+o(1))}{\ln^2 n}.
\]
The tail of the infinite series from $i=n$ onwards:
since $y^n = (n^{-1/(2n)})^n = n^{-1/2}$,
\begin{align*}
\sum_{i=n}^{\infty} i\,y^i
\;&=\; y^n\!\left(\frac{n}{1-y}+\frac{y}{(1-y)^2}\right)
\;=\; n^{-1/2}\cdot O\!\left(\frac{n^2}{\ln n}\right)
\\\;&=\; O\!\left(\frac{n^{3/2}}{\ln n}\right)
\;=\; o\!\left(\frac{n^2}{\ln^2 n}\right).
\end{align*}
Therefore $\sum_{i=1}^{n-1} i\,y^{i} = \Theta(n^2/\ln^2 n)$.
Since $n^{1/2+1/(2n)} = n^{1/2}\cdot n^{1/(2n)} = n^{1/2}(1+o(1))$,
\begin{equation}\label{eq:S-expoHD}
S \;=\; \Theta\!\left(\frac{n^{5/2}}{\log^2 n}\right).
\end{equation}

Finally, we follow the same arguments in Theorem~\ref{expohdOnLO} to establish the upper and the lower bound using $S$.
%
%\textbf{Upper bound.}
%Each improvement advances the \textsc{LeadingOnes} value by at least one,
%so the total expected number of evaluations satisfies
%$\mathbb{E}[T] \le \sum_{i=0}^{n-1} E(f_i) = S = O(n^{5/2}/\log^2 n)$.
%
%\textbf{Lower bound.}
%With probability $\Pr[x_1=0]=1/2$ the initial solution has
%$\textsc{LeadingOnes}$ value~$0$, so every level $i\in\{0,\dots,n-1\}$
%must be traversed.  Following the proof of Theorem~\ref{th:linHD-LO},
%the expected number of free-riders per improvement is at most~$2$,
%so each level~$i$ is visited with probability $p_i^*=\Omega(1)$.
%For the lower bound on~$E(f_i)$, we use the concentration
%$M \ge (1-o(1))\,n^{h_i}$ (w.o.p.), giving
%$E(f_i) \ge (1-o(1))\,i\cdot n^{1/2-(i-1)/(2n)} + n - i$.
%Hence
%\[
%\mathbb{E}[T]
%\;\ge\;
%\tfrac{1}{2}\sum_{i=0}^{n-1} p_i^*\cdot E(f_i)
%\;=\;
%\Omega(1)\cdot S
%\;=\;
%\Omega\!\left(\frac{n^{5/2}}{\log^2 n}\right).
%\]
\end{proof}


% ============================================================
%  DOWNSTREAM EDIT 1 — Speed-up commentary
%
%  The paragraph between % ============================================================
%  DROP-IN REPLACEMENT for Theorem \label{expohdOnLO}
%  (expoF on LeadingOnes)
%
%  Replaces:  \begin{theorem} \label{expohdOnLO} ...
%             through the matching \end{proof}
%
%  Also update the commentary sentence immediately after
%  \end{proof} — see note at bottom of this file.
% ============================================================

\begin{theorem} \label{expohdOnLO}
The {\oneoneIA } using \IPHfcm{} with {\expoF } optimises
\textsc{LeadingOnes} in $\Theta(n^3/\log^2 n)$ expected fitness function
evaluations.
\end{theorem}

\begin{proof}
At level~$i$ (i.e., with $i$ leading $1$-bits), the first bit flipped by the
hypermutation is chosen uniformly from $\{1,\dots,n\}$.  With probability
$i/n$ it hits a leading $1$-bit, destroying the prefix and using the full
potential of $n^{(n-i)/n}$ evaluations with no improvement; with probability
$1/n$ it hits the leftmost $0$-bit, yielding an improvement at cost~$1$; with
probability $(n-i-1)/n$ it hits a tail bit, costing $1$~evaluation with no
improvement.  By the law of total expectation,
\[
E(f_i)
= \frac{i}{n}\bigl(n^{(n-i)/n}+E(f_i)\bigr)
  +\frac{1}{n}\cdot 1
  +\frac{n-i-1}{n}\bigl(1+E(f_i)\bigr),
\]
and solving gives
\begin{equation}\label{eq:Efi-expoF}
E(f_i) \;=\; i\cdot n^{(n-i)/n} + n - i.
\end{equation}

We evaluate the sum
$S := \sum_{i=0}^{n-1} E(f_i)$, which governs both bounds.
Setting $x:=n^{-1/n}<1$,
\begin{align*}
S
&= \sum_{i=0}^{n-1}\bigl(i\cdot n^{(n-i)/n}+n-i\bigr)
 = n\sum_{i=1}^{n-1} i\,x^{i}
   \;+\;\underbrace{\sum_{i=0}^{n-1}(n-i)}_{\Theta(n^2)}.
\end{align*}
By the standard identity $\sum_{i=1}^{\infty} i\,x^i = x/(1-x)^2$
and Lemma~\ref{lem:expo} (with $\epsilon=1$),
\[
1-x \;=\; 1 - n^{-1/n}
    \;=\; \frac{\ln n}{n}\bigl(1+o(1)\bigr),
\]
so
\[
\frac{x}{(1-x)^2}
\;=\; \frac{1+o(1)}{\bigl(\frac{\ln n}{n}\bigr)^2}
\;=\; \frac{n^2(1+o(1))}{\ln^2 n}.
\]
The tail of the infinite series from $i=n$ onwards contributes
$O(n/\ln n)$ (since $x^n = n^{-1}$), which is $o(n^2/\ln^2 n)$.
Therefore $\sum_{i=1}^{n-1} i\,x^{i} = \Theta(n^2/\ln^2 n)$, giving
\begin{equation}\label{eq:S-expoF}
S \;=\; \Theta\!\left(\frac{n^3}{\log^2 n}\right).
\end{equation}

\textbf{Upper bound.}
Each improvement advances the \textsc{LeadingOnes} value by at least one,
so the total expected number of evaluations satisfies
$\mathbb{E}[T] \le \sum_{i=0}^{n-1} E(f_i) = S = O(n^3/\log^2 n)$.

\textbf{Lower bound.}
With probability $\Pr[x_1=0]=1/2$ the initial solution has
$\textsc{LeadingOnes}$ value~$0$, so every level $i\in\{0,\dots,n-1\}$
must be traversed.  Following the proof of Theorem~\ref{th:linHD-LO},
the expected number of free-riders per improvement is at most~$2$,
so each level~$i$ is visited with probability $p_i^*=\Omega(1)$.
Hence
\[
\mathbb{E}[T]
\;\ge\;
\tfrac{1}{2}\sum_{i=0}^{n-1} p_i^*\cdot E(f_i)
\;=\;
\Omega(1)\cdot S
\;=\;
\Omega\!\left(\frac{n^3}{\log^2 n}\right).
\]
\end{proof}

% ============================================================
%  DOWNSTREAM EDIT
%
%  The sentence immediately following \end{proof} currently reads:
%
%    "Theorem \ref{expohdOnLO} shows that the exponential based
%     mutation potential provides at least a logarithmic and at
%     most a $\sqrt{n}$ factor speed-up compared to static
%     mutation potentials."
%
%  Replace with:
%
%    "Theorem \ref{expohdOnLO} shows that the exponential-based
%     mutation potential provides a $\Theta(\log^2 n)$ factor
%     speed-up compared to static mutation potentials."
%
% ============================================================
 and the paragraph
%  introducing the expoHD theorem currently contains the sentence:
%
%    "An advantage of Hamming distance-based exponential decays ..."
%
%  This paragraph is fine as-is.  However, the sentence AFTER
%  \end{proof} of Theorem \ref{thm:expohdLO} that currently reads:
%
%    "We conclude that {\expoHD} provides larger hill-climbing
%     speed-ups compared to the other mutation potentials and is
%     stable to the scaling of fitness functions."
%
%  is still correct and needs no change.
%
%  IF there is a sentence earlier in the section reading something
%  like "at least a logarithmic and at most a $\sqrt{n}$ factor
%  speed-up", it should be updated to:
%
%    "Theorem~\ref{thm:expohdLO} shows that the Hamming-distance-based
%     exponential mutation potential provides a
%     $\Theta(\sqrt{n}\log^2 n)$ factor speed-up compared to
%     static mutation potentials."
%
%  (Since static = $\Theta(n^3)$ and expoHD = $\Theta(n^{5/2}/\log^2 n)$,
%   the speedup is $n^3/(n^{5/2}/\log^2 n) = \sqrt{n}\log^2 n$.)
% ============================================================


% ============================================================
%  DOWNSTREAM EDIT 2 — TABLE 2 CHANGES
%
%  In the table environment (label: table:afl), update:
%
%  OLD (expoHD row, LeadingOnes column):
%    $O(\frac{n^{5/2+\epsilon}}{\ln n})$, $\Omega(n^{9/4-\epsilon})$
%
%  NEW:
%    $\Theta(n^{5/2}/\log^2 n)$
% ============================================================
